% !TeX spellcheck = en_US
\documentclass[a4paper,10pt]{article}

\usepackage{xcolor}


\newcommand{\red}[1]{\textcolor{red}{#1}}

\usepackage{a4wide}


\begin{document}
%\pagenumbering{arabic}
\title{\#267 On Higher Order Program Verification via the Weighted Relational Model of Linear Logic\\ \emph{Detailed Answers}}
%\author{Beniamino Accattoli \and Ugo Dal Lago \and Gabriele Vanoni}
\maketitle

Points raised by the reviewers:

\begin{itemize}
	\item Review A:

\item Review B:
	\begin{itemize}
	\item \red{What do you mean by the point-wise partial order?}
	
	We indeed consider the pointwise order over the coefficients of the power series. 
	In Theorem 6.7, that the minimal solution of the system is indeed $a_G(1)$ can be shown using the fact that fixpoints in QRel are defined by iteration: given a polynomial equation $x=p(x,z)$, its lfp $a(z)$ is the sup of $a^0(z)=0$ and $a^{n+1}(z)=p(a^n(z),z)$, while
	the lfp y of $x=p(x,1)$ is the sup of $y^0=0, y^{n+1}=p(y^n,1)$; as evaluation in 1 is a Scott continous semiring homomorphism, we can prove by induction over n that $a^{n}(1)=y^n$ and then $a(1)=y$.
\end{itemize}


\item Review C:


\end{itemize}


\end{document}


